% praeambel.tex — Hamilton-Jacobi-kompendiet
% Den sædvanlige opsætning (jf. lineær algebra- og kvantemekanik-kompendierne)

\usepackage[a4paper,twoside,margin=2.5cm]{geometry}
\usepackage[danish]{babel}
\usepackage[utf8]{inputenc}
\usepackage[T1]{fontenc}
\usepackage{mathpazo}
\usepackage{amsmath,amssymb,amsthm}
\usepackage{physics}
\usepackage{graphicx}
\usepackage{tikz}
\usetikzlibrary{arrows.meta,calc,positioning,decorations.pathmorphing}
\usepackage{mdframed}
\usepackage{fancyhdr}
\usepackage{hyperref}
\usepackage{caption}
\captionsetup{font=small,margin=1.2cm}

% --- Sidehoved/-fod (even/odd) ---
\pagestyle{fancy}
\fancyhf{}
\fancyhead[LE,RO]{\thepage}
\fancyhead[RE]{\leftmark}
\fancyhead[LO]{\rightmark}
\renewcommand{\headrulewidth}{0.4pt}

% --- Farvede miljøer ---
\newmdenv[linecolor=blue!60!black,backgroundcolor=blue!5,linewidth=1pt,roundcorner=3pt]{eksempel}
\newmdenv[linecolor=orange!70!black,backgroundcolor=orange!5,linewidth=1pt,roundcorner=3pt]{bemaerkning}
\newmdenv[linecolor=orange!70!black,backgroundcolor=orange!5,linewidth=1pt,roundcorner=3pt]{opgave}

% --- Faste makroer (videreføres fra tidligere kompendier, udvides efter behov) ---
% \newcommand{\pd}[2]{\frac{\partial #1}{\partial #2}}
% Hamilton-Jacobi-specifikke makroer tilføjes her efterhånden som de opstår,
% f.eks. \newcommand{\Ham}{\mathcal{H}} for Hamilton-funktionen,
% \newcommand{\Lag}{\mathcal{L}} for Lagrange-funktionen,
% \newcommand{\Hact}{S} for virkningsfunktionen (Hamiltons principfunktion).
